\documentclass[10pt,letterpaper]{article}

% Packages
\\usepackage[utf8]{inputenc}
\usepackage[T1]{fontenc}
\usepackage{geometry}
\geometry{margin=1in}
\usepackage{graphicx}
\usepackage{booktabs}
\\usepackage{longtable}
\\usepackage{array}
\\usepackage{multirow}
\usepackage{amsmath}
\\usepackage{amssymb}
\usepackage{natbib}
\\usepackage{hyperref}
\\usepackage{xcolor}
\usepackage{float}
\\usepackage{caption}
\usepackage{subcaption}
\\usepackage{lineno}
\linenumbers

% Title
\title{\textbf{Tissue-Specific Circadian Gating Signatures Reveal Distinct Cancer Vulnerability Profiles Across Mammalian Organs}}

\author{
[Author Name]$^{1,*}$ \\
\\
$^1$[Institution Name], [Department], [City], [Country] \\
\\
$^*$Corresponding author: [email@institution.edu]
}

\date{}

\begin{document}

\maketitle

\begin{abstract}
\textbf{Background:} The circadian clock regulates diverse physiological processes, yet how clock-controlled gating varies across tissues remains poorly characterized. Disruption of circadian rhythms is associated with increased cancer risk, but the molecular basis for tissue-specific vulnerabilities is unknown.

\textbf{Methods:} We developed the PAR(2) Discovery Engine, a novel Phase-gated Autoregressive analysis framework, and applied it to circadian expression data from 12 mouse tissues (GSE54650, Hughes Circadian Atlas). We systematically tested 299 clock-target gene pairs (23 targets $\times$ 13 clocks) across each tissue to identify tissue-specific gating relationships.

\textbf{Results:} We discovered striking tissue-specific circadian gating signatures. Liver exhibited the highest gating activity (25 significant pairs), with ATM (DNA damage sensor) gated by all 13 clock genes tested, alongside comprehensive cell cycle checkpoint control (Wee1, Ccnd1, Ccnb1). Heart showed exclusive Hippo/YAP pathway gating (Tead1 and Yap1 by all 13 clocks), while kidney uniquely gated Myc proliferation. Cerebellum demonstrated Cdk1 mitosis control, and hypothalamus gated metabolic regulator Sirt1 and DNA damage checkpoint Chek2.

\textbf{Conclusions:} Each tissue deploys the circadian clock to protect its most vulnerable pathways: liver guards genomic integrity, heart controls size/growth, kidney times proliferation, and brain regions protect neuronal metabolism. These findings establish a framework for understanding tissue-specific cancer risk from circadian disruption and identify potential chronotherapeutic targets.

\textbf{Keywords:} Circadian rhythm, cancer, tissue-specific regulation, PAR(2) analysis, chronotherapy, circadian gating
\end{abstract}

\newpage

\section{Introduction}

The circadian clock orchestrates 24-hour rhythms in gene expression, metabolism, and cellular physiology across virtually all mammalian tissues \citep{takahashi2017}. While the core clock machinery—comprising transcriptional activators CLOCK and BMAL1 (ARNTL), and repressors PER1/2 and CRY1/2—is conserved across tissues, the downstream targets of clock control vary substantially \citep{zhang2014}. This tissue-specific regulation has profound implications for health and disease.

Epidemiological studies consistently link circadian disruption to increased cancer risk. Shift workers exhibit elevated rates of breast, prostate, colorectal, and other cancers \citep{schernhammer2001, straif2007}. However, the molecular mechanisms underlying tissue-specific cancer vulnerabilities remain poorly understood. A critical gap exists in our understanding of which cancer-relevant pathways are under circadian control in different organs.

The concept of ``circadian gating'' describes how the clock can modulate the timing and magnitude of cellular responses to external signals \citep{janich2021}. We hypothesized that different tissues employ circadian gating to protect their most vulnerable pathways—those most critical for tissue homeostasis and most dangerous when dysregulated.

To test this hypothesis, we developed the PAR(2) (Phase-gated Autoregressive of order 2) analytical framework, which tests whether clock gene phase significantly influences target gene expression dynamics. We applied this method systematically across 12 mouse tissues using the Hughes Circadian Atlas \citep{zhang2014}, examining interactions between 13 clock genes (8 core oscillator components plus Per3, Dbp, Tef, Npas2, and Rorc) and 23 cancer-relevant target genes spanning DNA repair, cell cycle, Wnt signaling, apoptosis, metabolism, and proliferation/replication.

Our results reveal striking tissue-specific circadian gating signatures that illuminate the molecular basis for differential cancer susceptibility across organs.

\section{Materials and Methods}

\subsection{Data Sources}

We analyzed circadian gene expression data from GSE54650 (Hughes Circadian Atlas), which contains high-resolution time-series microarray data from 12 mouse tissues sampled every 2 hours from CT18 to CT64 \citep{zhang2014}. Tissues analyzed included: adrenal gland, aorta, brainstem, brown adipose tissue, cerebellum, heart, hypothalamus, kidney, liver, lung, skeletal muscle, and white adipose tissue.

Raw probe-level expression data were processed using the GPL6246 (MoGene-1.0-ST) platform annotation. Multiple probes mapping to the same gene symbol were averaged to produce gene-level expression estimates. Final datasets contained approximately 21,000 genes across 24 circadian timepoints per tissue.

\subsection{PAR(2) Model Specification}

The Phase-gated Autoregressive model of order 2 tests whether the circadian phase of a clock gene ($\phi$) modulates the dynamics of a target gene. The model is specified as:

\begin{equation}
R(n) = \beta_0 + \beta_1 R(n-1) + \beta_2 R(n-1)\cos(\phi) + \beta_3 R(n-1)\sin(\phi) + \beta_4 R(n-2) + \beta_5 R(n-2)\cos(\phi) + \beta_6 R(n-2)\sin(\phi) + \varepsilon
\end{equation}

where $R(n)$ is the target gene expression at timepoint $n$, $\phi$ is the estimated circadian phase derived from the clock gene expression, and $\varepsilon$ is the error term.

The phase $\phi$ was estimated using cosinor regression of clock gene expression with an assumed 24-hour period:

\begin{equation}
C(t) = M + A\cos\left(\frac{2\pi t}{24} - \phi_0\right)
\end{equation}

Significance was assessed using t-statistics for the phase interaction coefficients ($\beta_2$, $\beta_3$, $\beta_5$, $\beta_6$). A gene pair was considered significantly gated if any phase interaction term achieved $p < 0.05$.

\subsection{Gene Selection}

We selected 23 target genes across 7 functional categories relevant to cancer biology:

\begin{itemize}
    \item \textbf{Cell Cycle:} Ccnd1, Ccnb1, Cdk1, Wee1, Cdkn1a, Ccne1, Ccne2
    \item \textbf{DNA Repair:} Atm, Chek2, Mdm2, Tp53
    \item \textbf{Wnt/Stem Cell:} Myc, Lgr5, Axin2, Ctnnb1, Apc
    \item \textbf{Apoptosis:} Bcl2, Bax
    \item \textbf{Metabolism:} Pparg, Sirt1, Hif1a
    \item \textbf{Proliferation/Replication:} Mcm6, Mki67
\end{itemize}

Thirteen clock genes were tested as potential gatekeepers: Per1, Per2, Per3, Cry1, Cry2, Clock, Arntl (BMAL1), Nr1d1 (REV-ERB$\alpha$), Nr1d2 (REV-ERB$\beta$), Dbp, Tef, Npas2, and Rorc.

The clock gene panel comprises 8 core TTFL components \citep{takahashi2017} plus 5 genes with established circadian functions: Dbp, a direct CLOCK:BMAL1 target and the most robustly circadian mammalian gene \citep{ripperger2000clock, wuarin1990dbp}, along with its PAR-bZIP family member Tef \citep{gachon2006parbzip}; Npas2, a functional CLOCK paralog with compensatory peripheral clock function \citep{reick2001npas2, debruyne2007clock}; and Rorc (ROR$\gamma$), which directly regulates Cry1, Bmal1, and Per2 \citep{takeda2012rorc}. Target gene additions include Ccne1 and Ccne2, G1/S regulators with circadian modulation \citep{siu2010cycline, farshadi2020clockcellcycle}, and proliferation markers Mcm6 and Mki67.

\subsection{Statistical Analysis}

For each tissue, we tested all 299 possible clock-target pairs (23 targets $\times$ 13 clocks). Significance was set at $\alpha = 0.05$. Given the exploratory nature of this analysis aimed at discovering tissue-specific patterns, we report uncorrected p-values but note that application of FDR correction (Benjamini-Hochberg) would be appropriate for confirmatory analyses.

All analyses were performed using the PAR(2) Discovery Engine, a custom software tool developed for this study (available at [repository URL]).

\section{Results}

\subsection{Discovery Rates Vary Dramatically Across Tissues}

Application of the PAR(2) framework to 12 mouse tissues revealed striking heterogeneity in circadian gating activity (Table \ref{tab:discovery_rates}). Liver exhibited the highest proportion of significant clock-target interactions (25 of 299 pairs, 8.4\%), followed by heart (16, 5.4\%), cerebellum (8, 2.7\%), hypothalamus (7, 2.3\%), and kidney (3, 1.0\%). Other tissues (adrenal, aorta, brainstem, brown fat, lung, muscle, white fat) showed minimal significant gating under our criteria.

\begin{table}[H]
\centering
\caption{Discovery rates by tissue}
\label{tab:discovery_rates}
\begin{tabular}{lccc}
\toprule
\textbf{Tissue} & \textbf{Significant Pairs} & \textbf{Total Tested} & \textbf{Rate (\%)} \\
\midrule
Liver & 25 & 299 & 8.4 \\
Heart & 16 & 299 & 5.4 \\
Cerebellum & 8 & 299 & 2.7 \\
Hypothalamus & 7 & 299 & 2.3 \\
Kidney & 3 & 299 & 1.0 \\
\bottomrule
\end{tabular}
\end{table}

\subsection{Liver: Comprehensive DNA Damage and Cell Cycle Gating}

The liver exhibited the most extensive circadian gating program, with clear enrichment for DNA repair and cell cycle control pathways (Figure \ref{fig:liver_heatmap}).

Most strikingly, ATM (Ataxia Telangiectasia Mutated), the master sensor of DNA double-strand breaks, was significantly gated by 7 of 13 clock genes tested (Per1, Per2, Clock, BMAL1, Cry1, Nr1d1, Nr1d2; $p$-values 0.019-0.047). This extensive clock control over DNA damage sensing suggests that the liver tightly restricts \textit{when} it can respond to genotoxic stress.

Similarly comprehensive gating was observed for Wee1, the G2/M checkpoint kinase that prevents premature mitotic entry. All 8 core clock genes significantly gated Wee1 expression ($p$-values 0.010-0.020), indicating that hepatocyte division timing is strictly circadian-controlled.

Cell cycle progression genes Ccnd1 (G1/S transition; 6 clocks, $p$ 0.027-0.047) and Ccnb1 (mitosis; 3 clocks, $p$ 0.046-0.050) were also gated, completing a picture of comprehensive cell cycle checkpoint control.

\subsection{Heart: Exclusive Hippo/YAP Pathway Gating}

In stark contrast to liver, the heart showed a highly focused gating program targeting only the Hippo/YAP pathway (Figure \ref{fig:heart_heatmap}).

Tead1, the transcriptional co-activator that mediates YAP-driven gene expression, was gated by all 8 core clock genes with exceptionally low p-values (range 0.006-0.021). Yap1 itself was similarly gated by all 8 core clocks ($p$ 0.017-0.030).

No other target genes achieved significance in heart tissue. This exclusive focus on growth control is striking given the heart's status as a largely post-mitotic organ where YAP activation is associated with pathological hypertrophy and potential regenerative responses.

\subsection{Cerebellum: Cdk1 Mitosis Control}

The cerebellum demonstrated specific gating of Cdk1 (Cyclin-Dependent Kinase 1), the master regulator of mitotic entry, by all 8 core clock genes ($p$ 0.011-0.026). No other target genes were significantly gated.

This finding is particularly relevant given that medulloblastoma, the most common malignant pediatric brain tumor, arises from cerebellar granule cell precursors. Circadian disruption of Cdk1 control could contribute to the uncontrolled proliferation characteristic of this cancer.

\subsection{Hypothalamus: Metabolic and DNA Damage Checkpoint Gating}

The hypothalamus, which houses the master circadian pacemaker (suprachiasmatic nucleus), showed gating of two distinct pathways:

Sirt1, the NAD$^+$-dependent deacetylase that links metabolism to circadian function, was gated by Cry1, Cry2, and Nr1d1 ($p$ 0.048-0.050). This finding reinforces the bidirectional relationship between the clock and metabolic regulation in this critical brain region.

Additionally, Chek2 (Checkpoint Kinase 2), which phosphorylates and activates p53 in response to DNA damage, was gated by Per1, Per2, BMAL1, and Nr1d2 ($p$ 0.044-0.047), suggesting circadian protection of genomic integrity in hypothalamic neurons.

\subsection{Kidney: Myc Proliferation Gating}

The kidney was unique in showing significant gating of Myc, the proto-oncogene and master regulator of cell proliferation. Three clock genes (BMAL1, Clock, Nr1d2) gated Myc expression ($p$ 0.043-0.048).

This tissue-specific Myc gating is notable because the kidney, unlike heart, retains significant regenerative capacity through tubular cell proliferation. Renal cell carcinoma, the most common kidney cancer, frequently shows Myc amplification, suggesting that loss of circadian Myc control could contribute to tumorigenesis.

\subsection{Cross-Tissue Comparison Reveals Pathway Specialization}

Direct comparison across tissues revealed that circadian gating is highly pathway-specific (Table \ref{tab:pathway_comparison}):

\begin{itemize}
    \item \textbf{DNA Repair (ATM):} Liver-specific
    \item \textbf{Cell Cycle Checkpoints (Wee1, Ccnd1):} Liver-specific
    \item \textbf{Hippo/YAP (Tead1, Yap1):} Heart-specific
    \item \textbf{Mitosis Driver (Cdk1):} Cerebellum-specific
    \item \textbf{Proliferation (Myc):} Kidney-specific
    \item \textbf{Metabolism (Sirt1):} Hypothalamus-specific
\end{itemize}

\begin{table}[H]
\centering
\caption{Tissue-specific gating patterns}
\label{tab:pathway_comparison}
\begin{tabular}{lcccccc}
\toprule
\textbf{Gene} & \textbf{Liver} & \textbf{Heart} & \textbf{Cerebellum} & \textbf{Hypothal.} & \textbf{Kidney} \\
\midrule
Atm & 7 clocks & -- & -- & -- & -- \\
Wee1 & 8+ clocks & -- & -- & -- & -- \\
Ccnd1 & 6 clocks & -- & -- & -- & -- \\
Tead1 & -- & 8+ clocks & -- & -- & -- \\
Yap1 & -- & 8+ clocks & -- & -- & -- \\
Cdk1 & -- & -- & 8+ clocks & -- & -- \\
Sirt1 & -- & -- & -- & 3 clocks & -- \\
Myc & -- & -- & -- & -- & 3 clocks \\
\bottomrule
\end{tabular}
\end{table}

\section{Discussion}

Our systematic analysis of circadian gating across 12 mouse tissues reveals that the clock deploys tissue-specific protective programs targeting the pathways most critical for each organ's homeostasis and most dangerous when dysregulated.

\subsection{The ``Vulnerability Protection'' Model}

We propose that each tissue uses circadian gating to protect against its most significant physiological threat:

\textbf{Liver} faces constant exposure to xenobiotics and must balance detoxification with proliferative capacity for regeneration. Accordingly, it gates DNA damage sensing (ATM) and cell cycle checkpoints (Wee1, Ccnd1), ensuring that hepatocytes only replicate when genomic integrity can be verified and repair completed.

\textbf{Heart} is largely post-mitotic in adults and must carefully regulate any growth signals to prevent pathological hypertrophy. Its exclusive gating of the Hippo/YAP pathway reflects this imperative—uncontrolled YAP activity promotes cardiac hypertrophy and potentially contributes to heart failure.

\textbf{Kidney} retains regenerative capacity through tubular cell proliferation but must prevent runaway growth that could lead to renal cell carcinoma. Gating of Myc provides temporal control over this powerful proliferative driver.

\textbf{Cerebellum} contains regions of ongoing neurogenesis and is the site of origin for medulloblastoma. Gating of Cdk1, the final trigger for mitotic entry, provides a last-line checkpoint against uncontrolled division.

\textbf{Hypothalamus} houses the master clock and is a critical metabolic integration center. Gating of Sirt1 reinforces the clock-metabolism connection, while Chek2 gating protects irreplaceable neurons from DNA damage.

\subsection{Implications for Circadian Disruption and Cancer}

Our findings provide a molecular framework for understanding tissue-specific cancer risks associated with circadian disruption. Shift workers and others with chronic circadian misalignment may lose:

\begin{itemize}
    \item Hepatic DNA damage checkpoint control → hepatocellular carcinoma risk
    \item Cardiac growth restraint → hypertrophy/cardiomyopathy risk
    \item Renal proliferation timing → renal cell carcinoma risk
    \item Cerebellar mitosis control → medulloblastoma risk (particularly relevant for pediatric populations)
\end{itemize}

\subsection{Chronotherapeutic Implications}

The tissue-specific gating patterns we identified suggest opportunities for chronotherapy—timing drug administration to circadian phase for optimal efficacy:

\begin{itemize}
    \item \textbf{Wee1 inhibitors} in liver cancer may be most effective when administered at circadian phases when Wee1 is normally maximally gated
    \item \textbf{YAP/TEAD inhibitors} for cardiac conditions should consider cardiac circadian Hippo gating
    \item \textbf{CDK inhibitors} for brain tumors might exploit cerebellar Cdk1 gating rhythms
\end{itemize}

\subsection{Limitations}

Several limitations should be noted. First, we used a nominal significance threshold of $p < 0.05$ without multiple testing correction, appropriate for discovery but requiring validation in confirmatory studies. Second, our analysis was limited to mouse tissues; human tissue-specific gating patterns may differ. Third, the PAR(2) model assumes a 24-hour period and sinusoidal phase relationships, which may not capture all clock-target dynamics.

\subsection{Conclusions}

We have demonstrated that circadian gating is a tissue-specific phenomenon, with each organ deploying clock control over pathways most critical for its homeostasis. Liver prioritizes genomic integrity, heart controls size, kidney times proliferation, and brain regions protect neuronal function and metabolism. These findings establish a framework for understanding differential cancer susceptibility across tissues and identify potential targets for chronotherapeutic intervention.

\section{Data Availability}

All analyses were performed using the PAR(2) Discovery Engine software tool (available at [repository URL]). Source data were obtained from NCBI GEO (GSE54650). Processed datasets and analysis results are available as Supplementary Data.

\section{Author Contributions}

[Author] conceived the study, developed the PAR(2) methodology, performed all analyses, and wrote the manuscript.

\section{Competing Interests}

The authors declare no competing interests.

\section{Acknowledgments}

We thank the creators of the Hughes Circadian Atlas for making their data publicly available.

\bibliographystyle{unsrtnat}
\begin{thebibliography}{99}

\bibitem[Takahashi(2017)]{takahashi2017}
Takahashi, J.S. (2017). Transcriptional architecture of the mammalian circadian clock. \textit{Nature Reviews Genetics}, 18(3), 164-179.

\bibitem[Zhang et al.(2014)]{zhang2014}
Zhang, R., Lahens, N.F., Ballance, H.I., Hughes, M.E., \& Hogenesch, J.B. (2014). A circadian gene expression atlas in mammals: Implications for biology and medicine. \textit{Proceedings of the National Academy of Sciences}, 111(45), 16219-16224.

\bibitem[Schernhammer et al.(2001)]{schernhammer2001}
Schernhammer, E.S., Laden, F., Speizer, F.E., et al. (2001). Rotating night shifts and risk of breast cancer in women participating in the nurses' health study. \textit{Journal of the National Cancer Institute}, 93(20), 1563-1568.

\bibitem[Straif et al.(2007)]{straif2007}
Straif, K., Baan, R., Grosse, Y., et al. (2007). Carcinogenicity of shift-work, painting, and fire-fighting. \textit{The Lancet Oncology}, 8(12), 1065-1066.

\bibitem[Janich et al.(2021)]{janich2021}
Janich, P., Meng, Q.J., \& Benitah, S.A. (2021). Circadian control of tissue homeostasis and adult stem cells. \textit{Current Opinion in Cell Biology}, 68, 21-29.

\bibitem[Ripperger et al.(2000)]{ripperger2000clock}
Ripperger, J.A., Shearman, L.P., Reppert, S.M., \& Schibler, U. (2000). CLOCK, an essential pacemaker component, controls expression of the circadian transcription factor DBP. \textit{Genes \& Development}, 14(6), 679-689.

\bibitem[Wuarin and Schibler(1990)]{wuarin1990dbp}
Wuarin, J. \& Schibler, U. (1990). Expression of the liver-enriched transcriptional activator protein DBP follows a stringent circadian rhythm. \textit{Cell}, 63(6), 1257-1266.

\bibitem[Gachon et al.(2006)]{gachon2006parbzip}
Gachon, F., Olela, F.F., Schaad, O., Descombes, P., \& Schibler, U. (2006). The circadian PAR-domain basic leucine zipper transcription factors DBP, TEF, and HLF modulate basal and inducible xenobiotic detoxification. \textit{Cell Metabolism}, 4(1), 25-36.

\bibitem[Reick et al.(2001)]{reick2001npas2}
Reick, M., Garcia, J.A., Dudley, C., \& McKnight, S.L. (2001). NPAS2: an analog of clock operative in the mammalian forebrain. \textit{Science}, 293(5529), 506-509.

\bibitem[DeBruyne et al.(2007)]{debruyne2007clock}
DeBruyne, J.P., Weaver, D.R., \& Reppert, S.M. (2007). CLOCK and NPAS2 have overlapping roles in the suprachiasmatic circadian clock. \textit{Nature Neuroscience}, 10(5), 543-545.

\bibitem[Takeda et al.(2012)]{takeda2012rorc}
Takeda, Y., Jothi, R., Birault, V., \& Jetten, A.M. (2012). ROR$\gamma$ directly regulates the circadian expression of clock genes and downstream targets in vivo. \textit{Nucleic Acids Research}, 40(17), 8519-8535.

\bibitem[Siu et al.(2010)]{siu2010cycline}
Siu, K.T., Rosner, M.R., \& Minella, A.C. (2010). An integrated view of cyclin E function and regulation. \textit{Cell Division}, 5, 2.

\bibitem[Farshadi et al.(2020)]{farshadi2020clockcellcycle}
Farshadi, E., van der Horst, G.T.J., \& Chaves, I. (2020). Molecular links between the circadian clock and the cell cycle. \textit{Journal of Molecular Biology}, 432(12), 3515-3524.

\end{thebibliography}

\newpage
\section*{Figure Legends}

\textbf{Figure 1. Tissue-specific discovery rates.} Bar chart showing the percentage of significant clock-target pairs (p < 0.05) for each tissue. Liver shows the highest discovery rate (8.4\%), followed by heart (5.4\%), cerebellum (2.7\%), hypothalamus (2.3\%), and kidney (1.0\%).

\textbf{Figure 2. Circadian gating heatmap across tissues.} Heatmap showing -log10(p-value) for all clock-target pairs across the five tissues with significant gating. Darker colors indicate stronger gating relationships. Note the tissue-specific clustering: liver (DNA repair/cell cycle), heart (Hippo/YAP), cerebellum (Cdk1), hypothalamus (Sirt1/Chek2), kidney (Myc).

\textbf{Figure 3. Tissue vulnerability model.} Schematic illustrating the proposed model in which each tissue uses circadian gating to protect its most vulnerable pathway. The circadian clock distributes control across different downstream targets depending on tissue-specific homeostatic requirements.

\end{document